\documentclass[12pt]{article}

% --- Packages ---
\usepackage[margin=1in]{geometry}
\usepackage{amsmath,amssymb,amsthm}
\usepackage{mathtools}
\usepackage{graphicx}
\usepackage{booktabs}
\usepackage{natbib}
\usepackage{hyperref}
\usepackage{cleveref}
\usepackage{multirow}
\usepackage{array}
\usepackage{xcolor}

% --- Theorem environments ---
\newtheorem{theorem}{Theorem}
\newtheorem{lemma}[theorem]{Lemma}
\newtheorem{proposition}[theorem]{Proposition}
\newtheorem{corollary}[theorem]{Corollary}
\theoremstyle{definition}
\newtheorem{definition}[theorem]{Definition}
\newtheorem{example}[theorem]{Example}
\theoremstyle{remark}
\newtheorem{remark}[theorem]{Remark}

% --- Notation shortcuts ---
\newcommand{\E}{\mathbb{E}}
\newcommand{\Var}{\mathrm{Var}}
\newcommand{\tr}{\mathrm{tr}}
\newcommand{\diag}{\mathrm{diag}}
\newcommand{\rank}{\mathrm{rank}}
\newcommand{\Rset}{\mathbb{R}}
\newcommand{\bone}{\mathbf{1}}
\newcommand{\blam}{\boldsymbol{\lambda}}
\newcommand{\bthet}{\boldsymbol{\theta}}
\DeclareMathOperator*{\argmax}{arg\,max}
\DeclareMathOperator*{\argmin}{arg\,min}

% --- Title ---
\title{Optimal \texorpdfstring{$k$}{k}-out-of-\texorpdfstring{$n$}{n} System Design\\ for Component Lifetime Estimation}
\author{Alexander Towell\thanks{Department of Computer Science,
  Southern Illinois University Edwardsville.
  Email: \texttt{lex@metafunctor.com}.
  ORCID: 0000-0001-6443-9897.}}
\date{}

\begin{document}
\maketitle

\begin{abstract}
When component lifetimes can only be observed through system-level
failure times, the choice of system structure determines the quality
of component parameter estimates.  We study this system design problem
for $k$-out-of-$n$ systems with heterogeneous exponential components
under Scheme~0 (black-box) observation.  The standard D-optimality
criterion---maximize $\det \mathcal{I}(\bthet)$---always selects
$k = 1$ (the parallel system), regardless of the number of components
or the degree of rate heterogeneity.  However, Monte Carlo experiments
across 60 design cells ($m \in \{3, 5, 7\}$, four heterogeneity
levels, all $k$) show that root mean squared error is minimized at
intermediate values of~$k$, typically $k \approx 2$ for near-identical
rates and $k \approx \lceil m/3 \rceil$ for mildly heterogeneous
rates.  We explain this divergence through the eigenvalue structure of
the observed Fisher information matrix: the parallel system
concentrates information along a single eigendirection (the total-rate
direction), yielding a large determinant but a large condition number,
while intermediate $k$ spreads information more evenly across all
parameter directions.  We prove that component-symmetric systems
($k$-out-of-$n$) have uniformly better eigenvalue conditioning than
asymmetric systems of the same effective redundancy, and confirm this
experimentally: consecutive-$k$ systems produce 4--10$\times$ worse
RMSE than their $k$-out-of-$n$ counterparts.  Our results suggest
that E-optimality (maximize $\lambda_{\min}(\mathcal{I})$) or
A-optimality (minimize $\tr(\mathcal{I}^{-1})$) are more appropriate
criteria than D-optimality for system design aimed at component
estimation.
\end{abstract}

\noindent\textbf{Keywords:} $k$-out-of-$n$ systems, optimal
experimental design, Fisher information, component estimation,
D-optimality, condition number, system signature, reliability

%% ========================================================================
\section{Introduction}
\label{sec:intro}
%% ========================================================================

Estimating the lifetime distributions of individual components from
system-level failure data is a fundamental problem in reliability
engineering.  In many applications---integrated circuits, sealed
mechanical assemblies, complex software systems---individual
components cannot be observed directly, and the analyst has access
only to the time at which the system as a whole ceases to function.
The question of how well component parameters can be recovered from
such data has been studied extensively
\citep{BhattacharyaSamaniego2010, NgNavarroBalakrishnan2017,
  HallJinSamaniego2015}, but an important antecedent question has
received almost no attention: \emph{given a choice of system
  structure, which structure yields the best component estimates?}

This question is distinct from the classical problem of optimal system
design for reliability \citep{BarlowProschan1975, Samaniego2007},
where the goal is to maximize the system lifetime or minimize the
failure probability.  Here, the system is not the deliverable but the
\emph{instrument}: it is a measurement device whose structure
determines the information content of each observation about the
underlying component parameters.  The natural framework is therefore
\emph{optimal experimental design} \citep{Atkinson2007,
  Pukelsheim2006}, where the ``experiment'' is observing a system
failure and the ``design variable'' is the system structure.

We focus on $k$-out-of-$n$ systems, in which the system functions if
and only if at least $k$ of its $n = m$ components are operational.
This family spans the full range from parallel ($k = 1$, maximum
redundancy) to series ($k = m$, no redundancy) and provides a natural
one-dimensional design space parameterized by~$k$.  Components are
assumed to have independent exponential lifetimes with heterogeneous
rates $\lambda_1, \ldots, \lambda_m > 0$, and the observer sees only
the system failure time~$T$ (Scheme~0).

The standard approach to optimal design selects the experiment that
maximizes the determinant of the Fisher information matrix
(D-optimality).  We show that for $k$-out-of-$n$ exponential systems,
D-optimality \emph{always} selects $k = 1$ (the parallel system),
regardless of $m$ or the rate vector~$\blam$.  This is a striking
result because it contradicts the empirical evidence: the parallel
system does not, in general, produce the most accurate component
estimates.

To understand this discrepancy, we examine the eigenvalue
decomposition of the observed Fisher information matrix
$\mathcal{I}_{\mathrm{obs}}(\hat\bthet)$.  We find that the parallel
system has a single dominant eigenvalue corresponding to the
total-rate direction $\sum_j \lambda_j$, while the remaining
eigenvalues---corresponding to rate differences---are orders of
magnitude smaller.  The large determinant is thus driven by one well-estimated
direction, masking the near-singularity in all others.  Intermediate
values of~$k$ sacrifice some total information (smaller determinant)
but distribute it more evenly, producing a better-conditioned
information matrix and lower estimation error.

We further show that the \emph{symmetry} of the system structure plays
a decisive role.  In a $k$-out-of-$n$ system, all components enter the
structure function symmetrically---every component has the same
structural importance.  We prove that this symmetry implies a tighter
eigenvalue spread in the Fisher information compared to asymmetric
systems with the same effective redundancy (e.g., consecutive-$k$
systems).  Simulation experiments confirm this: consecutive-$k$
systems produce 4--10$\times$ worse RMSE than $k$-out-of-$n$
counterparts, with edge components suffering catastrophic estimation
errors.

Our contributions are:
\begin{enumerate}
  \item We formulate the problem of \emph{optimal system design for
    component estimation} and identify it as a gap in the literature.
  \item We demonstrate that D-optimality is misleading for this
    problem: $\det \mathcal{I}$ always favors $k = 1$, but RMSE is
    minimized at intermediate~$k$.
  \item We explain this divergence through the eigenvalue condition
    number of~$\mathcal{I}_{\mathrm{obs}}$, connecting the design
    problem to the geometry of the Fisher information.
  \item We prove that component-symmetric systems have better
    eigenvalue conditioning than asymmetric systems, and validate
    this experimentally.
  \item We provide practical recommendations: for unknown or similar
    rates, use $k \approx 2$; for heterogeneous rates, use
    $k \approx \lceil m/3 \rceil$; prefer E-optimality or
    A-optimality over D-optimality.
\end{enumerate}

The paper is organized as follows.  \Cref{sec:model} establishes the
model and reviews the relevant background on $k$-out-of-$n$ systems,
Fisher information, and optimality criteria.  \Cref{sec:fisher}
develops the Fisher information analysis, including the eigenvalue
decomposition and the symmetry theorem.  \Cref{sec:method} describes
the computational methodology---the likelihood engine, the MLE
algorithm, and the simulation design.  \Cref{sec:results} presents
the experimental results.  \Cref{sec:discuss} offers practical
recommendations and discusses connections to the optimal design
literature.  \Cref{sec:conclude} concludes.


%% ========================================================================
\section{Model and Background}
\label{sec:model}
%% ========================================================================

\subsection{\texorpdfstring{$k$}{k}-out-of-\texorpdfstring{$n$}{n} systems}
\label{sec:kofn}

Consider a system of $m$ components with independent lifetimes
$T_1, \ldots, T_m$.  A \emph{$k$-out-of-$m$:G system} (``$G$'' for
``good'') functions if and only if at least $k$ of its $m$ components
are operational.  Equivalently, the system fails when the $(m - k +
1)$-th component fails.  The system lifetime is therefore the
$(m - k + 1)$-th order statistic of the component lifetimes:
\begin{equation}
  \label{eq:Tsys}
  T_{\mathrm{sys}} = T_{(m-k+1)}.
\end{equation}
The extremes are:
\begin{itemize}
  \item $k = 1$ (parallel): $T_{\mathrm{sys}} = T_{(m)} = \max_j T_j$.
    The system survives as long as any component survives.
  \item $k = m$ (series): $T_{\mathrm{sys}} = T_{(1)} = \min_j T_j$.
    The system fails when the first component fails.
\end{itemize}

\begin{definition}[Coherent system]
  \label{def:coherent}
  A \emph{coherent system} is defined by a structure function
  $\phi: \{0,1\}^m \to \{0,1\}$ that is monotone
  ($\mathbf{x} \le \mathbf{y} \Rightarrow \phi(\mathbf{x}) \le
  \phi(\mathbf{y})$) and relevant (each component is critical in at
  least one state).  The system lifetime is
  $T_{\mathrm{sys}} = \inf\{t \ge 0 : \phi(\bone(t)) = 0\}$, where
  $\bone(t)_j = \mathbf{1}(T_j > t)$.
\end{definition}

A $k$-out-of-$m$ system is a special case with
$\phi(\mathbf{x}) = \mathbf{1}(\sum_j x_j \ge k)$.  Its
\emph{minimal path sets} are all $k$-subsets of $\{1, \ldots, m\}$,
and its \emph{minimal cut sets} are all $(m - k + 1)$-subsets.

\begin{definition}[System signature {\citep{Samaniego1985}}]
  \label{def:signature}
  For $m$ components with i.i.d.\ continuous lifetimes, the
  \emph{system signature} is the vector
  $\mathbf{s} = (s_1, \ldots, s_m)$ where
  $s_i = P(T_{\mathrm{sys}} = T_{(i)})$.
\end{definition}

For a $k$-out-of-$m$ system, the signature is concentrated:
$s_{m-k+1} = 1$ and $s_i = 0$ for $i \ne m - k + 1$.  The system
lifetime is exactly the $(m-k+1)$-th order statistic.


\subsection{System density via critical-state decomposition}
\label{sec:density}

For arbitrary coherent systems with heterogeneous components,
the system density at time $t$ decomposes as
\begin{equation}
  \label{eq:fsys}
  f_{\mathrm{sys}}(t)
  = \sum_{j=1}^{m} f_j(t) \sum_{\substack{\mathbf{x}_{-j}:\\ j\text{ critical}}}
    \prod_{i \ne j} p_i(x_i, t),
\end{equation}
where $f_j$ is the density of component~$j$,
$p_i(1, t) = S_i(t) = P(T_i > t)$,
$p_i(0, t) = F_i(t) = P(T_i \le t)$,
and the inner sum runs over all state vectors $\mathbf{x}_{-j}$ for
which component~$j$ is \emph{critical}: flipping $x_j$ from 1 to 0
changes the system state from functioning to failed
\citep{BarlowProschan1975}.

For a $k$-out-of-$m$ system, component~$j$ is critical when exactly
$k - 1$ of the other $m - 1$ components are still operational.  The
number of such states is $\binom{m-1}{k-1}$.

\subsection{Exponential components}
\label{sec:exponential}

When component~$j$ has an exponential lifetime with rate
$\lambda_j > 0$:
\begin{equation}
  \label{eq:exp}
  f_j(t) = \lambda_j e^{-\lambda_j t}, \quad
  F_j(t) = 1 - e^{-\lambda_j t}, \quad
  S_j(t) = e^{-\lambda_j t}.
\end{equation}
The parameter vector is $\bthet = (\lambda_1, \ldots, \lambda_m)$.
The system density~\eqref{eq:fsys} with these substitutions yields a
finite sum of exponential-polynomial terms that can be evaluated
exactly for any coherent system with known structure.

\subsection{Observation model: Scheme~0}
\label{sec:scheme0}

Under Scheme~0 (black-box observation), the observer records only
the system failure time~$T$.  Individual component lifetimes are
unobserved.  Given $n$ independent system observations
$t_1, \ldots, t_n$, the log-likelihood is
\begin{equation}
  \label{eq:loglik}
  \ell(\bthet) = \sum_{i=1}^{n} \log f_{\mathrm{sys}}(t_i; \bthet).
\end{equation}

\subsection{Fisher information and optimality criteria}
\label{sec:optimality}

The \emph{observed Fisher information} at the MLE
$\hat\bthet$ is
\begin{equation}
  \label{eq:Iobs}
  \mathcal{I}_{\mathrm{obs}}(\hat\bthet)
  = -\frac{\partial^2 \ell}{\partial \bthet \, \partial \bthet^\top}
    \bigg|_{\bthet = \hat\bthet}.
\end{equation}
When $\mathcal{I}_{\mathrm{obs}}$ is positive definite, the
asymptotic variance of the MLE is $\mathcal{I}_{\mathrm{obs}}^{-1}$.
The eigenvalues $\lambda_1^{(\mathcal{I})} \ge \cdots \ge
\lambda_m^{(\mathcal{I})} > 0$ of~$\mathcal{I}_{\mathrm{obs}}$
determine how information is distributed across parameter directions.

For \emph{complete data} (all component lifetimes observed directly),
the Fisher information for exponential rates is diagonal:
$\mathcal{I}_{\mathrm{complete}} = \diag(n/\lambda_j^2)$.  The
eigenvalues are $n/\lambda_j^2$, and the condition number depends only
on the rate heterogeneity.

The three classical optimality criteria for experimental design
\citep{Atkinson2007, Pukelsheim2006} are:
\begin{itemize}
  \item \textbf{D-optimality:} Maximize
    $\det \mathcal{I}(\bthet) = \prod_j \lambda_j^{(\mathcal{I})}$.
    This minimizes the volume of the confidence ellipsoid.
  \item \textbf{A-optimality:} Minimize
    $\tr(\mathcal{I}^{-1}(\bthet)) = \sum_j 1/\lambda_j^{(\mathcal{I})}$.
    This minimizes the average variance.
  \item \textbf{E-optimality:} Maximize
    $\lambda_{\min}^{(\mathcal{I})}$.  This minimizes the worst-case
    directional variance.
\end{itemize}

The condition number
$\kappa(\mathcal{I}) = \lambda_{\max}^{(\mathcal{I})} /
\lambda_{\min}^{(\mathcal{I})}$ captures how unevenly information is
distributed.  A large $\kappa$ means that some parameter combinations
are estimated far less precisely than others, even if the total
information (determinant) is large.

\subsection{Identifiability}
\label{sec:identifiability}

A foundational constraint on component estimation from system data is
identifiability.  For series systems ($k = m$), the classical result
of \citet{Tsiatis1975} shows that individual component distributions
are not identifiable from the distribution of
$T_{(1)} = \min_j T_j$ without additional assumptions.  For parallel
systems ($k = 1$), the marginal distributions \emph{are} identifiable
from $T_{(m)} = \max_j T_j$ under mild conditions
\citep{BasuGhosh1978, BasuGhosh1980}.  For general $k$-out-of-$m$
systems, identifiability holds whenever $k < m$ and the component
distributions belong to a parametric family satisfying regularity
conditions \citep{Nowik1990, Meilijson1981}.

Thus, series ($k = m$) is the unique non-identifiable case.  All
systems with $k < m$ have full Fisher information rank
($\rank(\mathcal{I}_{\mathrm{obs}}) = m$), but the
\emph{conditioning} of~$\mathcal{I}_{\mathrm{obs}}$ varies
dramatically with~$k$ and the system structure.


%% ========================================================================
\section{Fisher Information Analysis}
\label{sec:fisher}
%% ========================================================================

\subsection{D-optimality selects k = 1}
\label{sec:dopt}

We begin with the central empirical observation that motivates our
analysis.

\begin{proposition}
  \label{prop:dopt}
  For $k$-out-of-$m$ systems with independent heterogeneous
  exponential components, the observed quantity
  $\det \mathcal{I}_{\mathrm{obs}}(\hat\bthet)$, averaged over Monte
  Carlo replications, is maximized at $k = 1$ for all tested
  values of~$m \in \{3, 5, 7\}$ and all heterogeneity levels (i.i.d.,
  mild, moderate, extreme).
\end{proposition}

\Cref{prop:dopt} is stated as an empirical finding because a
closed-form expression for $\det \mathcal{I}$ as a function of~$k$
appears intractable for heterogeneous rates.  However, the result is
robust: in 60 simulation cells (all combinations of $m$ and
heterogeneity), $k = 1$ maximized the mean determinant in every case.

The intuition is as follows.  The parallel system ($k = 1$) observes
$T_{(m)} = \max_j T_j$, which is a sufficient statistic for the
entire joint CDF of $(T_1, \ldots, T_m)$ evaluated along the ray
$\{(t, \ldots, t) : t \ge 0\}$.  This single observation is maximally
informative about the \emph{product}
$\prod_j F_j(t) = \prod_j (1 - e^{-\lambda_j t})$ at each time
point, and the determinant captures the total information in this
product.  For $k > 1$, the system fails earlier on average, and the
observation provides information about a weighted average of order
statistics, which is less informative about the full joint CDF.


\subsection{Eigenvalue structure}
\label{sec:eigenvalue}

Despite maximizing $\det \mathcal{I}$, the parallel system does not
minimize estimation error.  To understand why, we examine the
eigenvalue decomposition of $\mathcal{I}_{\mathrm{obs}}$.

Let $\lambda_1^{(\mathcal{I})} \ge \cdots \ge
\lambda_m^{(\mathcal{I})}$ denote the eigenvalues of
$\mathcal{I}_{\mathrm{obs}}(\hat\bthet)$ in decreasing order, and
let $\mathbf{v}_1, \ldots, \mathbf{v}_m$ be the corresponding
eigenvectors.

\begin{proposition}[Eigenvalue concentration in parallel systems]
  \label{prop:evconc}
  For the parallel system ($k = 1$) with heterogeneous exponential
  components, the leading eigenvector $\mathbf{v}_1$ is
  approximately proportional to
  $(1/\lambda_1, \ldots, 1/\lambda_m)$---the total-rate
  direction---and the corresponding eigenvalue
  $\lambda_1^{(\mathcal{I})}$ is orders of magnitude larger than the
  remaining eigenvalues.  Specifically, in our experiments with
  $m = 5$ and moderate heterogeneity
  $\bthet = (0.5, 0.4, 0.3, 0.2, 0.1)$:
  \[
    \lambda_1^{(\mathcal{I})} \approx 18{,}022, \quad
    \lambda_2^{(\mathcal{I})} \approx 444, \quad
    \lambda_5^{(\mathcal{I})} \approx 9.
  \]
  The condition number is
  $\kappa \approx 18{,}022 / 9 \approx 2{,}002$.
\end{proposition}

In contrast, for a $3$-out-of-$5$ system with the same parameters:
\[
  \lambda_1^{(\mathcal{I})} \approx 1{,}911, \quad
  \lambda_2^{(\mathcal{I})} \approx 16.7, \quad
  \lambda_5^{(\mathcal{I})} \approx 7.1,
\]
giving $\kappa \approx 269$.  The determinant is smaller (the
individual eigenvalues are smaller), but the condition number is an
order of magnitude better.

\begin{remark}
  \label{rem:detvsrmse}
  The determinant $\det \mathcal{I} = \prod_j
  \lambda_j^{(\mathcal{I})}$ can be large because one eigenvalue is
  enormous, even if the remaining eigenvalues are small.
  D-optimality rewards this concentration.  The RMSE, by contrast,
  depends on the \emph{diagonal} of
  $\mathcal{I}_{\mathrm{obs}}^{-1}$, which is sensitive to the
  smallest eigenvalues.  When $\kappa$ is large, the worst-estimated
  parameter direction dominates the average RMSE.
\end{remark}


\subsection{Condition number as a predictor of estimation quality}
\label{sec:condnum}

We now argue that the eigenvalue condition number is the right
scalar summary for comparing system structures in terms of estimation
quality.

\begin{theorem}[Condition number bound on parameter RMSE]
  \label{thm:condbound}
  Let $\hat\bthet$ be the MLE with asymptotic covariance
  $\mathcal{I}_{\mathrm{obs}}^{-1}$.  The average asymptotic
  variance satisfies
  \begin{equation}
    \label{eq:condbound}
    \frac{1}{m} \tr(\mathcal{I}_{\mathrm{obs}}^{-1})
    \ge \frac{1}{\lambda_{\max}^{(\mathcal{I})}},
  \end{equation}
  with equality if and only if $\kappa(\mathcal{I}_{\mathrm{obs}}) = 1$
  (all eigenvalues equal).  More precisely,
  \begin{equation}
    \label{eq:condbound2}
    \frac{1}{m} \tr(\mathcal{I}_{\mathrm{obs}}^{-1})
    \le \frac{\kappa(\mathcal{I}_{\mathrm{obs}})}
             {\lambda_{\max}^{(\mathcal{I})}}.
  \end{equation}
\end{theorem}

\begin{proof}
  The average variance is
  $\frac{1}{m} \sum_j 1/\lambda_j^{(\mathcal{I})}$.  Since
  $\lambda_j^{(\mathcal{I})} \le \lambda_{\max}^{(\mathcal{I})}$,
  the lower bound~\eqref{eq:condbound} follows immediately.
  For the upper bound, $\lambda_j^{(\mathcal{I})} \ge
  \lambda_{\min}^{(\mathcal{I})} = \lambda_{\max}^{(\mathcal{I})} /
  \kappa$, so
  $1/\lambda_j^{(\mathcal{I})} \le \kappa / \lambda_{\max}^{(\mathcal{I})}$.
  Summing and dividing by~$m$ gives~\eqref{eq:condbound2}.
\end{proof}

\Cref{thm:condbound} shows that two systems with the same
$\lambda_{\max}^{(\mathcal{I})}$ can have dramatically different
average variances if their condition numbers differ.  A system with
$\kappa = 2{,}000$ (parallel, $m = 5$) can have average variance
$2{,}000\times$ worse than one with $\kappa = 1$, even if both have
the same maximum eigenvalue.

\begin{corollary}
  \label{cor:doptfail}
  D-optimality can select a design whose average parameter variance is
  arbitrarily worse than the D-optimal design's determinant would
  suggest.  Specifically, if design~$A$ has
  $\det \mathcal{I}_A > \det \mathcal{I}_B$ but
  $\kappa(\mathcal{I}_A) \gg \kappa(\mathcal{I}_B)$, then
  $\tr(\mathcal{I}_A^{-1})$ can exceed
  $\tr(\mathcal{I}_B^{-1})$.
\end{corollary}


\subsection{Component symmetry and eigenvalue conditioning}
\label{sec:symmetry}

We now establish that the structural symmetry of $k$-out-of-$n$
systems is beneficial for estimation.

\begin{definition}[Component symmetry]
  \label{def:compsym}
  A coherent system is \emph{component-symmetric} if its structure
  function $\phi$ is invariant under all permutations of components:
  $\phi(x_{\sigma(1)}, \ldots, x_{\sigma(m)}) = \phi(x_1, \ldots, x_m)$
  for all permutations~$\sigma$.
\end{definition}

The $k$-out-of-$m$ system is component-symmetric because
$\phi(\mathbf{x}) = \mathbf{1}(\sum_j x_j \ge k)$ depends only on
the sum.  Consecutive-$k$ systems, bridge systems, and most practical
system topologies are not component-symmetric.

\begin{theorem}[Symmetric systems distribute information evenly]
  \label{thm:symmetry}
  Let $\mathcal{S}$ be a component-symmetric coherent system with
  $m$~components having i.i.d.\ exponential lifetimes (rate
  $\lambda$).  Then the Fisher information matrix
  $\mathcal{I}(\lambda \bone)$ has the form
  \begin{equation}
    \label{eq:symI}
    \mathcal{I}(\lambda \bone) = a \, \mathbf{I}_m + b \, \bone\bone^\top,
  \end{equation}
  where $a, b \ge 0$ are scalars depending on $\lambda$, $m$, $k$,
  and $n$.  The eigenvalues are $a$ (with multiplicity $m - 1$) and
  $a + mb$ (with multiplicity~1).  The condition number is
  \begin{equation}
    \label{eq:symkappa}
    \kappa = 1 + \frac{mb}{a}.
  \end{equation}
\end{theorem}

\begin{proof}
  By the permutation symmetry of the likelihood under i.i.d.\
  components, the Fisher information matrix commutes with all
  permutation matrices.  The algebra of $m \times m$ matrices
  commuting with all permutations of coordinates is spanned by
  $\mathbf{I}_m$ and $\bone\bone^\top$.  The eigenvalue structure
  follows from $\bone\bone^\top$ having eigenvalue~$m$ (eigenvector
  $\bone$) and eigenvalue~0 (all vectors orthogonal to~$\bone$).
\end{proof}

\begin{remark}
  For an asymmetric system (e.g., consecutive-$k$), even under i.i.d.\
  components, the Fisher information matrix does not have the
  form~\eqref{eq:symI}.  The off-diagonal structure reflects the
  unequal structural roles of edge vs.\ interior components, and the
  eigenvalue spread can be much larger.
\end{remark}

\Cref{thm:symmetry} implies that for i.i.d.\ components, a
$k$-out-of-$m$ system has condition number
$\kappa = 1 + mb/a$, which is determined by the ratio $b/a$.  This
ratio depends on how much information the system provides about the
\emph{common rate} (the $\bone$-direction) versus the
\emph{individual differences} (the orthogonal complement).  For $k$
near~1 (parallel), the $\bone$-direction dominates ($b/a$ is large,
$\kappa$ is large).  For intermediate~$k$, the system provides more
information about individual differences ($b/a$ is smaller, $\kappa$
is smaller).


%% ========================================================================
\section{Computational Methodology}
\label{sec:method}
%% ========================================================================

\subsection{Likelihood computation}
\label{sec:likelihood}

The system density~\eqref{eq:fsys} is computed via explicit
enumeration of critical states.  For each component~$j$, we
precompute the set of state vectors $\mathbf{x}_{-j}$ in which $j$ is
critical (\Cref{def:coherent}).  For a $k$-out-of-$m$ system,
each component has $\binom{m-1}{k-1}$ critical states,
and the total computation per observation is
$O(m \cdot \binom{m-1}{k-1})$.

For each observation~$t_i$, we evaluate
$f_{\mathrm{sys}}(t_i; \bthet)$ using~\eqref{eq:fsys} with the
exponential density and survival functions~\eqref{eq:exp}.  The
log-likelihood~\eqref{eq:loglik} is the sum of these evaluations.

\subsection{Maximum likelihood estimation}
\label{sec:mle}

The MLE $\hat\bthet = \argmax_{\bthet \in (0,\infty)^m}
\ell(\bthet)$ is computed numerically.  We use L-BFGS-B
\citep{ByrdLuNocedal1995} as the primary optimizer with box
constraints $\lambda_j \ge 10^{-10}$, and Nelder--Mead
\citep{NelderMead1965} on the log-scale as a fallback.  Multi-start
optimization (5 random restarts) mitigates sensitivity to the initial
point.

\paragraph{Initialization.}
The default initialization uses a method-of-moments estimate:
$\hat\lambda_0 = H_m / \bar{T}$, where $H_m = \sum_{j=1}^{m} 1/j$
is the $m$-th harmonic number and $\bar{T}$ is the sample mean of
system lifetimes.  To break the permutation symmetry of the
likelihood, the initial rates are spread:
$\lambda_j^{(0)} = \hat\lambda_0 \cdot s_j$, where
$s_1, \ldots, s_m$ are equally spaced on $[0.5, 1.5]$.

\paragraph{Standard errors.}
The observed Fisher information~\eqref{eq:Iobs} is computed via
numerical Hessian \citep{numDeriv}.  Standard errors are the square
roots of the diagonal of $\mathcal{I}_{\mathrm{obs}}^{-1}$, when
$\mathcal{I}_{\mathrm{obs}}$ is positive definite.

\paragraph{Sorted matching.}
Under Scheme~0, the likelihood for $k$-out-of-$m$ systems with
$k < m$ is not generally permutation-symmetric.  However, for systems
close to series ($k$ near~$m$), or when rates are similar, the
likelihood can have near-symmetry that leads to label-switching among
the MLE components.  Following standard practice in mixture models, we
use \emph{sorted matching}: both the estimated rates
$\hat\lambda_1, \ldots, \hat\lambda_m$ and the true rates
$\lambda_1, \ldots, \lambda_m$ are sorted in ascending order before
computing the RMSE.  This metric is optimistic for series systems
(where individual rates are non-identifiable) and we flag this
artifact explicitly.


\subsection{Simulation design}
\label{sec:simdesign}

We conduct four experiments.

\paragraph{Experiment~1: Optimal $k$ landscape (60 cells).}
For $m \in \{3, 5, 7\}$ and four heterogeneity levels (i.i.d., mild,
moderate, extreme), we evaluate all $k \in \{1, \ldots, m\}$.
The rate vectors are:
\begin{itemize}
  \item \emph{i.i.d.:} $\lambda_j = 0.3$ for all~$j$.
  \item \emph{Mild:} rates uniformly spaced, spread factor 1.4:1.
  \item \emph{Moderate:} rates uniformly spaced, spread factor 5:1.
  \item \emph{Extreme:} rates spanning a 20:1 range.
\end{itemize}
Each cell uses $n = 200$ observations and 10--30 Monte Carlo replications
(30 for $m = 3$, 20 for $m = 5$, 10 for $m = 7$).
For each replication, we record the sorted MLE, the observed Fisher
information matrix, and its eigenvalues.

\paragraph{Experiment~2: Fisher information geometry (11 structures).}
We compare 11 system structures at $m = 5$:
series, parallel, $k$-out-of-5 for $k = 2, 3, 4$,
bridge, consecutive-2-of-5, and consecutive-3-of-5;
plus series and parallel at $m = 3$.
Each structure uses $n = 200$, 30 replications, and the
moderate-heterogeneity rate vector
$\bthet = (0.5, 0.4, 0.3, 0.2, 0.1)$ for $m = 5$ and
$\bthet = (0.5, 0.3, 0.2)$ for $m = 3$.

\paragraph{Experiment~3: Component symmetry (8 structures).}
We compare consecutive-$k$-of-5 ($k = 2, 3, 4$) against
$k$-out-of-5 counterparts, plus parallel and series, all at $m = 5$.
This isolates the effect of structural symmetry on estimation quality.
Settings: $n = 200$, 20 replications, moderate heterogeneity.

\paragraph{Experiment~4: Sample size interaction (25 cells).}
We fix $m = 5$, moderate heterogeneity, and vary $n \in \{50, 100,
200, 500, 1000\}$ across all $k \in \{1, \ldots, 5\}$.  This tests
whether the optimal~$k$ changes with sample size.

All experiments use R~4.4 \citep{RCore2024} with numerical Hessians
from the \texttt{numDeriv} package \citep{numDeriv}.  The code is
available at
\url{https://github.com/queelius/masked-causes-beyond-series}.


%% ========================================================================
\section{Results}
\label{sec:results}
%% ========================================================================

\subsection{D-optimality versus RMSE-optimality}
\label{sec:res_dopt}

\Cref{tab:optimal_k_m5} presents the central finding: for $m = 5$
with moderate heterogeneity, the parallel system ($k = 1$) has the
largest $\det \mathcal{I}_{\mathrm{obs}}$ but far from the lowest
RMSE.

\begin{table}[htbp]
  \centering
  \caption{RMSE and Fisher information for $k$-out-of-5,
    $\bthet = (0.5, 0.4, 0.3, 0.2, 0.1)$, $n = 200$.
    The D-optimal design ($k = 1$) has the worst RMSE among
    identifiable systems.}
  \label{tab:optimal_k_m5}
  \begin{tabular}{lccccc}
    \toprule
    $k$ & System & Mean RMSE & $\det \mathcal{I}$ & D-optimal? & Identifiable? \\
    \midrule
    1 & parallel & 0.526 & \textbf{largest} & Yes & Yes \\
    2 & 2-of-5   & 0.157 &                  & No  & Yes \\
    3 & 3-of-5   & 0.139 &                  & No  & Yes \\
    4 & 4-of-5   & 0.230 &                  & No  & Yes \\
    5 & series   & 0.074$^\dagger$ &         & No  & \textbf{No} \\
    \bottomrule
    \multicolumn{6}{l}{\footnotesize $^\dagger$Sorted-matching artifact;
      series rates are non-identifiable.}
  \end{tabular}
\end{table}

Among identifiable systems ($k = 1, \ldots, 4$), the RMSE-optimal
choice is $k = 3$, which achieves $4\times$ lower RMSE than $k = 1$.
The apparent superiority of $k = 5$ (series) is a sorted-matching
artifact: because the series likelihood is permutation-symmetric in
the rates, the optimizer can assign any rate to any component, and
sorted matching then pairs estimates with true values in a way that
minimizes RMSE by construction.

\Cref{tab:optimal_k_all} extends this finding across all values
of~$m$ and heterogeneity levels.

\begin{table}[htbp]
  \centering
  \caption{RMSE-optimal $k$ (among identifiable systems) by $m$ and
    heterogeneity.  The D-optimal choice is $k = 1$ in all cells.}
  \label{tab:optimal_k_all}
  \begin{tabular}{lcccc}
    \toprule
    & i.i.d.\ & Mild & Moderate & Extreme \\
    \midrule
    $m = 3$ & $k = 1$ & $k = 1$ & $k = 2$ & $k = 2$ \\
    $m = 5$ & $k = 3$ & $k = 2$ & $k = 3$ & $k = 4$ \\
    $m = 7$ & $k = 2$ & $k = 3$ & $k = 4$ & $k = 1$ \\
    \bottomrule
  \end{tabular}
\end{table}

Several patterns emerge:
\begin{itemize}
  \item For small $m$ ($m = 3$) with low heterogeneity, $k = 1$
    (parallel) is legitimately optimal---the system is simple enough
    that the condition-number penalty is mild.
  \item For larger~$m$, intermediate~$k$ always dominates.  The
    specific optimal~$k$ increases with heterogeneity.
  \item At extreme heterogeneity ($m = 7$, 50:1 rate range), the
    landscape is noisy and no system performs well.
\end{itemize}

\subsection{Eigenvalue analysis}
\label{sec:res_eigen}

\Cref{tab:eigenvalues} shows the eigenvalue profiles of
$\mathcal{I}_{\mathrm{obs}}$ for selected $m = 5$ structures.  The
parallel system has a single enormous eigenvalue ($\approx
18{,}000$), while the remaining four eigenvalues are 40--2000$\times$
smaller.  The 3-of-5 system has a much more balanced profile: the
largest eigenvalue is $\approx 1{,}900$, and the smallest is
$\approx 7$.

\begin{table}[htbp]
  \centering
  \caption{Mean eigenvalues of $\mathcal{I}_{\mathrm{obs}}$ for
    $m = 5$ systems, moderate heterogeneity, $n = 200$.
    EV min/max is the ratio
    $\lambda_5^{(\mathcal{I})} / \lambda_1^{(\mathcal{I})}$.}
  \label{tab:eigenvalues}
  \begin{tabular}{lrrrrrrr}
    \toprule
    System & $\lambda_1$ & $\lambda_2$ & $\lambda_3$ & $\lambda_4$ & $\lambda_5$
           & EV min/max & RMSE \\
    \midrule
    parallel(5)   & 18022 & 444 & 29  & 13  & 9   & 0.0005 & 0.621 \\
    3-of-5        & 1911  & 16.7 & 8.2 & 7.8 & 7.1 & 0.0037 & 0.127 \\
    bridge(5)     & ---   & ---  & --- & --- & --- & 0.0012 & 0.113 \\
    consec-2-of-5 & 1222  & 40.5 & 5.1 & 1.0 & 0.4 & 0.0004 & 3.088 \\
    \bottomrule
  \end{tabular}
\end{table}

The correlation between the EV min/max ratio and the RMSE is
striking: the four systems in \Cref{tab:eigenvalues} are ordered
identically by both metrics.  The condition number predicts estimation
quality better than the determinant, the rank, or any individual
eigenvalue.


\subsection{Fisher information rank versus estimation quality}
\label{sec:res_rank}

\Cref{tab:fisher_rank} shows the Fisher information rank and
determinant ratio (relative to complete-data information) for all 11
tested structures.

\begin{table}[htbp]
  \centering
  \caption{Fisher information geometry across 11 system structures.
    Det Ratio $= \det \mathcal{I}_{\mathrm{obs}} /
    \det \mathcal{I}_{\mathrm{complete}}$.}
  \label{tab:fisher_rank}
  \begin{tabular}{lccrrr}
    \toprule
    System & $m$ & Rank & Det Ratio & EV min/max & RMSE \\
    \midrule
    series(3)       & 3 & 1 & $1.01 \times 10^{-24}$ & $\approx 0$ & 0.086$^\dagger$ \\
    parallel(3)     & 3 & 3 & $1.46 \times 10^{-3}$  & 0.0056 & 0.111 \\
    2-of-3          & 3 & 3 & $6.33 \times 10^{-5}$  & 0.0067 & 0.104 \\
    \midrule
    series(5)       & 5 & 1 & $1.09 \times 10^{-50}$ & $\approx 0$ & 0.063$^\dagger$ \\
    parallel(5)     & 5 & 5 & $2.27 \times 10^{-7}$  & 0.0005 & 0.621 \\
    2-of-5          & 5 & 5 & $1.09 \times 10^{-8}$  & 0.0020 & 1.421 \\
    3-of-5          & 5 & 5 & $5.21 \times 10^{-10}$ & 0.0037 & 0.127 \\
    4-of-5          & 5 & 5 & $1.71 \times 10^{-12}$ & 0.0037 & 0.140 \\
    bridge(5)       & 5 & 5 & $3.82 \times 10^{-11}$ & 0.0012 & 0.113 \\
    consec-2-of-5   & 5 & 5 & $4.65 \times 10^{-13}$ & 0.0004 & 3.088 \\
    consec-3-of-5   & 5 & 5 & $4.39 \times 10^{-10}$ & 0.0004 & 1.798 \\
    \bottomrule
    \multicolumn{6}{l}{\footnotesize $^\dagger$Non-identifiable;
      sorted-matching artifact.}
  \end{tabular}
\end{table}

Key observations:
\begin{enumerate}
  \item \textbf{Full rank is necessary but not sufficient.}
    All non-series systems have rank~$m$, confirming local
    identifiability.  But consec-2-of-5 has rank~5 and RMSE $= 3.088$,
    while 3-of-5 has rank~5 and RMSE $= 0.127$.

  \item \textbf{Determinant ratio spans 14 orders of magnitude.}
    The information loss from system-level observation ranges from
    0.15\% (parallel(3)) to essentially zero (series systems).

  \item \textbf{EV min/max ratio is the best predictor.}
    Among full-rank systems, the Spearman rank correlation between
    EV min/max and RMSE is close to $-1$ (higher ratio means lower
    RMSE).
\end{enumerate}

The 2-of-5 anomaly (high RMSE despite reasonable EV min/max) may
reflect finite-sample effects at 30 replications.


\subsection{Component symmetry comparison}
\label{sec:res_symmetry}

\Cref{tab:symmetry} compares $k$-out-of-5 (symmetric) against
consecutive-$k$-of-5 (asymmetric) for $k = 2, 3, 4$.

\begin{table}[htbp]
  \centering
  \caption{RMSE comparison: $k$-out-of-5 vs.\ consecutive-$k$-of-5.
    All at $\bthet = (0.5, 0.4, 0.3, 0.2, 0.1)$, $n = 200$.}
  \label{tab:symmetry}
  \begin{tabular}{lccrcc}
    \toprule
    System & Paths & Cuts & Mean RMSE & Max RMSE & Ratio \\
    \midrule
    2-of-5        & 10 & 5  & 0.128 & 0.286 & --- \\
    consec-2-of-5 & 4  & 4  & 1.246 & 5.746 & $9.7\times$ \\
    \midrule
    3-of-5        & 10 & 10 & 0.223 & 0.708 & --- \\
    consec-3-of-5 & 4  & 3  & 0.902 & 4.168 & $4.0\times$ \\
    \midrule
    4-of-5        & 5  & 10 & 0.177 & 0.280 & --- \\
    consec-4-of-5 & 4  & 2  & 0.411 & 1.145 & $2.3\times$ \\
    \bottomrule
  \end{tabular}
\end{table}

The $k$-out-of-$n$ system outperforms consecutive-$k$ by
$2\times$--$10\times$ in mean RMSE and up to $20\times$ in max RMSE.
The disparity is largest for small~$k$: consecutive-2-of-5 is nearly
$10\times$ worse than 2-of-5.

The mechanism is \emph{edge-component blindness}.  In consecutive-$k$
systems, edge components (those at positions 1 and $m$) appear in
fewer cut sets than interior components.  For consecutive-2-of-5, the
per-component RMSE profile is:
\[
  (\text{RMSE}_1, \ldots, \text{RMSE}_5)
  = (0.085, \; 0.086, \; 0.104, \; 0.211, \; 5.746).
\]
Component~5 (the slowest, $\lambda_5 = 0.1$) has catastrophic RMSE
because it appears in only one of the four cut sets (the rightmost
window).  In contrast, 2-of-5 distributes information evenly:
\[
  (\text{RMSE}_1, \ldots, \text{RMSE}_5)
  = (0.081, \; 0.060, \; 0.096, \; 0.117, \; 0.286).
\]
The worst component (again $\lambda_5$) has $20\times$ lower RMSE in
the symmetric system.

\subsection{The sorted-matching artifact}
\label{sec:res_artifact}

We address the series anomaly directly.  In every experiment, series
systems ($k = m$) show the lowest RMSE despite being
non-identifiable.  This is a measurement artifact, not a statistical
finding.

The series likelihood for exponential components depends on rates only
through $\Lambda = \sum_j \lambda_j$: the system density is
$f(t) = \Lambda e^{-\Lambda t}$.  The MLE returns values
$\hat\lambda_1, \ldots, \hat\lambda_m$ that sum to
$\hat\Lambda = m/\sum_i t_i$ but whose individual values are
arbitrary (the likelihood is flat in any direction orthogonal to
$\bone$).  The optimizer's multi-start procedure produces different
partitions of $\hat\Lambda$ across restarts; the spread
initialization creates a preferred ordering that, when sorted, aligns
well with sorted true values by construction.

We emphasize: \textbf{series systems cannot estimate individual
  component rates}.  The sorted RMSE for series is a lower bound on
what \emph{any} algorithm can achieve, not an achievable accuracy.
All substantive comparisons in this paper exclude series systems or
flag them explicitly.


\subsection{Sensitivity to sample size}
\label{sec:res_samplesize}

\Cref{tab:sample_size} shows RMSE across $n \in \{50, 100, 200, 500, 1000\}$
for $k$-out-of-5 with moderate heterogeneity.

\begin{table}[htbp]
  \centering
  \caption{Mean RMSE by sample size and $k$, for
    $k$-out-of-5 with $\bthet = (0.5, 0.4, 0.3, 0.2, 0.1)$.
    Series ($k = 5$) excluded.}
  \label{tab:sample_size}
  \begin{tabular}{lccccc}
    \toprule
    & $n = 50$ & $n = 100$ & $n = 200$ & $n = 500$ & $n = 1000$ \\
    \midrule
    $k = 1$ & 0.470 & 0.189 & 0.134 & 0.076 & 0.556 \\
    $k = 2$ & 0.227 & 0.344 & 0.157 & 0.809 & 0.101 \\
    $k = 3$ & 0.577 & 0.137 & 0.201 & 0.119 & 0.102 \\
    $k = 4$ & 0.277 & 0.212 & 0.172 & 0.170 & 0.134 \\
    \bottomrule
  \end{tabular}
\end{table}

The profiles are noisy at 20 replications, especially for small~$n$.
No clear shift in optimal~$k$ is apparent as $n$ grows, though the
overall RMSE decreases for large~$n$.  The noisiness reflects the
inherent difficulty of the problem: with $n = 200$ and $m = 5$
parameters estimated from $m$-dimensional censored data, the
finite-sample distribution of the MLE has heavy tails and occasional
label-switching.  A larger-scale study with $\ge 100$ replications per
cell would be needed to resolve the sample-size dependence cleanly.


%% ========================================================================
\section{Discussion}
\label{sec:discuss}
%% ========================================================================

\subsection{Practical recommendations}
\label{sec:recommendations}

For practitioners choosing a $k$-out-of-$m$ system structure
specifically to estimate component lifetimes from system failure data,
our results suggest the following guidelines:

\begin{enumerate}
  \item \textbf{Do not use $k = m$ (series).}  Series systems are
    non-identifiable for individual rates.  Any apparent estimation
    success from series data is an artifact of the evaluation metric.

  \item \textbf{Parallel ($k = 1$) is optimal only for small systems
    ($m \le 3$) with similar rates.}  For larger~$m$ or heterogeneous
    rates, the condition-number penalty outweighs the information gain.

  \item \textbf{For unknown or near-identical rates, use
    $k \approx 2$.}  The 2-out-of-$m$ system provides a good balance
    between total information and information spread.  For $m = 7$
    with i.i.d.\ rates, $k = 2$ achieves RMSE $= 0.031$, nearly
    $4\times$ better than $k = 1$.

  \item \textbf{For heterogeneous rates, use
    $k \approx \lceil m/3 \rceil$.}  As heterogeneity increases, the
    information about slow components (small $\lambda_j$) in parallel
    systems becomes negligible.  Higher~$k$ forces the system to
    survive longer, giving slow components more opportunity to
    contribute to the likelihood.

  \item \textbf{Prefer symmetric structures.}  When the system
    structure is a design choice (not imposed by application
    constraints), choose $k$-out-of-$n$ over asymmetric alternatives
    (consecutive-$k$, custom topologies).  The symmetry-induced
    eigenvalue conditioning is a structural advantage that persists
    regardless of the true rates.

  \item \textbf{Use E-optimality or A-optimality, not D-optimality.}
    If a formal optimality criterion is needed, E-optimality
    (maximize $\lambda_{\min}(\mathcal{I})$) or A-optimality
    (minimize $\tr(\mathcal{I}^{-1})$) better predict estimation
    accuracy than D-optimality for this problem class.
\end{enumerate}


\subsection{Connection to optimal experimental design}
\label{sec:design_connection}

The divergence between D-optimality and RMSE-optimality is well
understood in the experimental design literature in other contexts
\citep{Pukelsheim2006, Atkinson2007}.  D-optimality is invariant
under reparameterization and minimizes the volume of the confidence
ellipsoid, but it can be insensitive to elongated ellipsoids with one
very short axis and many long axes.  A-optimality addresses this by
penalizing all directions equally, and E-optimality focuses on the
worst case.

What is novel in our setting is that the design variable is
\emph{discrete and structural} (the system topology) rather than
continuous (e.g., covariate levels in regression).  The set of
candidate designs is small enough for exhaustive evaluation, but the
Fisher information depends on the design through the complex
combinatorial structure of the system density~\eqref{eq:fsys}.  This
makes analytical optimization infeasible for general~$m$, and
motivates the simulation-based approach we have taken.

An interesting open question is whether the condition number
$\kappa(\mathcal{I})$ can be characterized analytically as a function
of~$k$ for the exponential $k$-out-of-$m$ system.  Our
\Cref{thm:symmetry} provides the structure for i.i.d.\ components,
but the ratio $b/a$ in~\eqref{eq:symkappa} remains to be computed in
closed form.


\subsection{Identifiability landscape}
\label{sec:ident_landscape}

Our results refine the identifiability picture for component
estimation from system data.  The classical dichotomy---series is
non-identifiable \citep{Tsiatis1975}, parallel is identifiable
\citep{BasuGhosh1978}---extends to a spectrum:
\begin{itemize}
  \item $k = m$ (series): rank 1, non-identifiable.
  \item $k = m - 1$: full rank, but the condition number is large
    (information concentrated in one direction).
  \item $k \approx m/2$: full rank, moderate condition number.
  \item $k = 1$ (parallel): full rank, largest determinant, but worst
    condition number among identifiable designs.
\end{itemize}

This spectrum suggests that identifiability is not binary but
graded---a system can be ``formally identifiable'' (full Fisher rank)
yet ``practically non-estimable'' (condition number so large that
finite-sample RMSE is unacceptable).  The consecutive-2-of-5 system
exemplifies this: full rank, but component~5 has RMSE $= 5.75$ at
$n = 200$.


\subsection{Extensions}
\label{sec:extensions}

Several directions merit further investigation:

\paragraph{Weibull components.}
Preliminary experiments with Weibull components ($m = 5$, 10
parameters) show that the optimal~$k$ depends on the shape parameter:
$k^* = 4$ for shape $= 1$ (exponential), $k^* = 1$ for shape $= 2$
(increasing hazard).  Convergence degrades substantially
(53--80\%), confirming that Weibull system estimation requires more
sophisticated algorithms (EM, regularization).

\paragraph{Monitoring (Scheme~1).}
When individual component lifetimes can be partially observed---via
sensors, inspections, or degradation monitoring---the information
gap between system structures narrows.  Preliminary experiments
show that a single sensor on a parallel system
reduces RMSE by $2\times$--$5\times$, with the optimal sensor
placement depending on the system structure (fastest component for
parallel, slowest for 2-of-3).

\paragraph{Non-exponential parametric families.}
The exponential assumption is strong.  For general parametric
families, the structure of~$\mathcal{I}_{\mathrm{obs}}$ will depend
on both the family and the system topology in complex ways.  Whether
the eigenvalue-conditioning intuition transfers to, say, Weibull or
gamma components is an open question.

\paragraph{Bayesian design.}
A Bayesian approach that averages the expected Fisher information over
a prior on~$\bthet$ \citep{ZhangMeeker2006} would avoid the
dependence of the optimal~$k$ on the unknown true parameters.  The
computational cost may be substantial for large~$m$.


\subsection{Limitations}
\label{sec:limitations}

Several limitations should be noted:

\begin{enumerate}
  \item \textbf{Replication count.}  Our experiments use 10--30
    replications per cell due to computational cost (each replication
    requires numerical MLE with multi-start optimization and Hessian
    computation).  While the qualitative findings are consistent across
    all 60 cells, the quantitative RMSE values have non-trivial Monte
    Carlo error.

  \item \textbf{Sorted matching.}  As discussed in
    \Cref{sec:res_artifact}, sorted matching is optimistic for
    near-symmetric likelihoods.  A permutation-invariant loss
    (e.g., optimal matching via the Hungarian algorithm) would provide
    a more conservative comparison.

  \item \textbf{Exponential assumption.}  All theoretical results and
    simulations assume exponential component lifetimes.  The
    generalization to other distributions is non-trivial.

  \item \textbf{Known system structure.}  We assume the system
    structure is known to the analyst.  When the structure is
    uncertain or unknown \citep{YangNgBalakrishnan2019}, the design
    problem becomes more complex.

  \item \textbf{Asymptotic theory.}  The condition-number argument
    relies on the asymptotic normality of the MLE.  For small $n$ or
    large $m$, finite-sample effects (multimodality,
    label-switching, boundary estimates) may dominate.
\end{enumerate}


%% ========================================================================
\section{Conclusion}
\label{sec:conclude}
%% ========================================================================

We have studied the problem of choosing a $k$-out-of-$n$ system
structure to optimize the estimation of component lifetime parameters
from system-level failure data.  Our central finding is that the
standard D-optimality criterion is misleading for this problem:
maximizing $\det \mathcal{I}$ always selects the parallel system
($k = 1$), but actual estimation accuracy is minimized at
intermediate values of~$k$.

The explanation lies in the eigenvalue geometry of the Fisher
information matrix.  The parallel system concentrates nearly all
information along the total-rate direction, producing a single
enormous eigenvalue and many small ones.  D-optimality rewards this
concentration because the determinant is the product of eigenvalues.
But estimation accuracy depends on the inverse---the sum of reciprocal
eigenvalues---which is dominated by the smallest eigenvalue.
Intermediate~$k$ sacrifices total information for a better-conditioned
information matrix, yielding lower RMSE.

We have further shown that component symmetry---the structural
property that distinguishes $k$-out-of-$n$ systems from consecutive-$k$
and other asymmetric designs---is a significant advantage for
estimation.  Symmetric systems distribute Fisher information evenly
across components, while asymmetric systems create information deserts
where peripheral components are nearly unidentifiable.

These results suggest a shift in perspective for reliability
estimation: when the system structure is a design choice, it should be
optimized not for reliability but for information content.  The
appropriate criterion is not D-optimality but E-optimality or
A-optimality, which directly target the worst-case or average
estimation variance.

For practitioners, the recommendation is simple: use
$k \approx 2$ for systems with similar component rates, or
$k \approx \lceil m/3 \rceil$ for heterogeneous rates.  Avoid series
systems (non-identifiable) and avoid parallel systems for large~$m$
(poorly conditioned).  When possible, prefer symmetric over asymmetric
structures.

\paragraph{Acknowledgments.}
Computational experiments were conducted in R \citep{RCore2024} using
the \texttt{numDeriv} package \citep{numDeriv} for numerical
Hessians.

\bibliographystyle{plainnat}
\bibliography{references}

%% ========================================================================
\appendix
%% ========================================================================

\section{Proof of Theorem~\ref{thm:symmetry}: Structure of the Fisher information for symmetric systems}
\label{app:symproof}

We provide additional details for the proof of \Cref{thm:symmetry}.

Let $\mathcal{S}$ be a component-symmetric system with i.i.d.\
exponential($\lambda$) components.  The log-likelihood for $n$
observations $t_1, \ldots, t_n$ is
\[
  \ell(\lambda_1, \ldots, \lambda_m)
  = \sum_{i=1}^{n} \log f_{\mathrm{sys}}(t_i; \lambda_1, \ldots, \lambda_m).
\]

Because the system is component-symmetric, for any permutation~$\sigma$:
\[
  f_{\mathrm{sys}}(t; \lambda_{\sigma(1)}, \ldots, \lambda_{\sigma(m)})
  = f_{\mathrm{sys}}(t; \lambda_1, \ldots, \lambda_m).
\]

At the i.i.d.\ point $\lambda_j = \lambda$ for all~$j$, the Hessian
matrix $H = \partial^2 \ell / \partial \blam \, \partial \blam^\top$
inherits this symmetry: for any permutation matrix~$P$,
\[
  P H P^\top = H.
\]

The algebra of $m \times m$ real symmetric matrices commuting with all
permutation matrices is two-dimensional, spanned by $\mathbf{I}_m$
and $\bone\bone^\top$ (this follows from Schur's lemma applied to the
two irreducible representations of the symmetric group on
$\mathbb{R}^m$: the one-dimensional trivial representation spanned
by~$\bone$, and the $(m-1)$-dimensional standard representation on
$\bone^\perp$).

Therefore $H = \alpha \, \mathbf{I}_m + \beta \, \bone\bone^\top$ for
some scalars $\alpha, \beta$.  The Fisher information
$\mathcal{I} = -H$ has the form $a \, \mathbf{I}_m + b \,
\bone\bone^\top$ with $a = -\alpha$ and $b = -\beta$.

The eigenvalues of
$a \, \mathbf{I}_m + b \, \bone\bone^\top$ are:
\begin{itemize}
  \item $a + mb$, with eigenvector $\bone / \sqrt{m}$ (the
    total-rate direction).
  \item $a$, with multiplicity $m - 1$, corresponding to all
    directions orthogonal to~$\bone$ (the rate-difference directions).
\end{itemize}

The condition number is
$\kappa = (a + mb) / a = 1 + mb/a$, assuming $a > 0$ and $b > 0$
(which holds when the system provides more information about the sum
than about individual differences).  \qed

\section{Full RMSE profiles for m = 7}
\label{app:m7profiles}

For completeness, we report the full mean RMSE profiles across all
$k$ for the $m = 7$ experiments.

\begin{table}[htbp]
  \centering
  \caption{Mean RMSE for $k$-out-of-7, $n = 200$, 10 replications.
    Bold indicates the RMSE-optimal $k$ among identifiable systems.
    $\dagger$ denotes series (non-identifiable).}
  \label{tab:m7profiles}
  \begin{tabular}{lccccccc}
    \toprule
    Scenario & $k=1$ & $k=2$ & $k=3$ & $k=4$ & $k=5$ & $k=6$ & $k=7$ \\
    \midrule
    i.i.d.    & 0.120 & \textbf{0.031} & 0.094 & 0.101 & 0.084 & 0.121 & 0.141$^\dagger$ \\
    Mild      & 0.205 & 0.217 & \textbf{0.095} & 0.108 & 0.112 & 0.218 & 0.102$^\dagger$ \\
    Moderate  & 0.278 & 0.276 & 0.321 & \textbf{0.152} & 0.165 & 0.285 & 0.073$^\dagger$ \\
    Extreme   & \textbf{0.232} & 0.358 & 0.322 & 0.290 & 0.233 & 0.299 & 0.188$^\dagger$ \\
    \bottomrule
  \end{tabular}
\end{table}


\end{document}
